% =============================================================================
%   II. RELATED WORK
% =============================================================================
\section{Related Work}
\label{sec:related}

\subsection{ESG Rating Divergence and Measurement}

Commercial ESG ratings disagree substantially. Berg, K\"{o}lbel and Rigobon
\cite{berg_2022} attribute the divergence mainly to measurement and scope
differences rather than weighting. Chatterji et al.\ \cite{chatterji_2016}
document low convergence, and disagreement itself carries return and risk
implications \cite{gibson_2021, avramov_2022}. Christensen, Serafeim and
Sikochi \cite{christensen_2022} show that more disclosure can \emph{increase}
disagreement, and Kotsantonis and Serafeim \cite{kotsantonis_2019} catalogue
data problems such as inconsistent peer sets, imputation and scale. This
literature treats ESG ratings as noisy, model-dependent measurements. We apply
the same lens to a public-data ESG score, where every modelling choice is
visible and can be audited cell by cell.

\subsection{ESG, Returns and Risk}

Meta-analytic evidence reports predominantly non-negative ESG--performance
associations \cite{friede_2015, eccles_2014}. Materiality matters: industry-material
issues drive the effect \cite{khan_2016}. Equilibrium models imply that
investor tastes for sustainability can \emph{lower} the expected returns of
green assets while making them hedges \cite{pastor_2021, pedersen_2021}, and
carbon exposure is priced \cite{bolton_2021}. On the risk side, CSR is associated with lower
systematic risk \cite{albuquerque_2019}, and high-trust firms fared better in
the 2008--2009 crisis \cite{lins_2017}. A risk association alone does not
identify an ESG effect, however. Low-volatility and low-beta stocks
\cite{ang_2006, frazzini_pedersen_2014} and high-quality firms
\cite{asness_2019} have distinct risk and return profiles, and realised
volatility is highly persistent. An ESG score that loads on these
characteristics will ``predict'' risk without containing ESG information.
We therefore test the risk channel against these characteristics and against
recent volatility, and trace it back to the sources of the score's inputs.

\subsection{Public-Data and Proxy ESG}

Where direct disclosure is sparse, researchers infer ESG-relevant behaviour
from observables. Examples include environmental performance and firm value
\cite{konar2001does}, eco-efficiency \cite{derwall2005ecoefficiency,
guenster2011ecoefficiency}, employee satisfaction \cite{edmans_2011},
entrenchment indices \cite{bebchuk2009what}, machine-learned ESG scores
\cite{damato_2022} and text-mined sustainability reports
\cite{luccioni_2020}. Each proxy has a rationale, but aggregating many proxies
built from the \emph{same} financial statements risks producing an ESG score
that is a relabelled financial-quality score. We quantify this risk directly
(Section~\ref{sec:res_construct}).

\subsection{Multi-Factor Models and Composite Indicators}

Factor models explain the cross-section of returns with a small number of
characteristics \cite{fama_french_1993, carhart_1997, fama_french_2015,
asness_2013, jegadeesh_titman_1993}. Commercial ESG indices tilt or exclude
constituents on ESG ratings under tracking-error constraints \cite{giese_2019,
alighanbari_2022, msci_esg}. Our index is a multi-criteria \emph{ranking}
rather than a pricing model. Its construction follows the OECD handbook for
composite indicators \cite{oecd_2008}: normalisation, weighting, aggregation
and sensitivity analysis.

\subsection{Validation, Multiple Testing and Overfitting}

Hundreds of published factors fail to replicate once multiple testing is
taken into account \cite{harvey_2016}. Backtests that choose among many
configurations overfit with high probability \cite{bailey_2017}. The remedy
proposed across empirical sciences is to fix the hypotheses and the analysis
plan before outcomes are observed \cite{nosek_2018}. We adopt this remedy in
two steps. The first is a hold-out window with hypotheses fixed before
outcomes were computed and a family-wise error correction \cite{holm_1979}.
The second is a formally pre-registered window whose scores, plan and code
are hashed before it opens. We also show concretely how in-sample designs
common in index research produce favourable evidence irrespective of
predictive validity.
